\documentclass[a4paper,14pt]{extarticle}

\usepackage[T2A]{fontenc}
\usepackage[utf8]{inputenc}
\usepackage[russian]{babel}

\usepackage[left=25mm, right=15mm, top=20mm, bottom=20mm]{geometry}
\usepackage{graphicx}
\usepackage{hyperref}
\usepackage{xcolor}
\usepackage{indentfirst}
\usepackage{fvextra}

\definecolor{codegray}{gray}{0.95}

\DefineVerbatimEnvironment{code}{Verbatim}{
    breaklines=true,
    breakanywhere=true,
    frame=single,
    framerule=0.4pt,
    fontsize=\small,
    bgcolor=codegray,
    rulecolor=\color{black},
    samepage=false
}

\hypersetup{
    colorlinks=true,
    linkcolor=blue,
    urlcolor=blue
}

\begin{document}

\begin{titlepage}
    \centering
    {\scshape\large Министерство науки и высшего образования Российской Федерации\par}
    \vspace{0.5cm}
    {\scshape\large Федеральное государственное автономное образовательное учреждение\\
    высшего образования\par}
    {\scshape\large «Национальный исследовательский университет ИТМО»\par}
    \vspace{0.5cm}
    {\large Институт математики\par}
    \vspace{3cm}

    {\Large\bfseries Лабораторная работа №0 (вводная)\par}
    \vspace{0.5cm}
    {\large Знакомство с системой контроля версий git и сервисом GitHub\par}
    \vspace{3cm}

    \begin{flushright}
        \begin{tabular}{l}
            Студент: \textit{Бузырев Иван Михайлович}\\
            Группа: \textit{R-3341}\\
            Преподаватель: \textit{Шаврин А.А.}
        \end{tabular}
    \end{flushright}

    \vfill
    {\large Санкт-Петербург\\ \the\year\par}
\end{titlepage}

\tableofcontents
\newpage

\section{Текст задания}

\subsection*{Задание 1}
\textbf{Цель:} поразмышлять о своих целях и задачах на текущий семестр.

\begin{enumerate}
    \item С помощью веб-интерфейса GitHub создать файл \texttt{goals.md} в репозитории.
    \item Используя язык разметки markdown, создать заголовки <<Мой опыт>> и <<Мои цели>>.
    \item Написать один абзац (1--5 предложений) о своём опыте по разработке ПО. Добавить работающую ссылку (например, на профиль в GitHub).
    \item Создать нумерованный список из 2--3 целей на курс.
    \item Создать ненумерованный список из 3--9 задач для достижения целей.
    \item Создать папку \texttt{resources} в корне репозитория, загрузить в неё изображение (\texttt{.png/.jpg/.jpeg}) и добавить его в \texttt{.md} файл.
    \item Закоммитить файлы в репозиторий на GitHub.
    \item Проверить, что \texttt{.md} файл отформатирован и отображается корректно.
\end{enumerate}

\subsection*{Задание 2}
\textbf{Цель:} научиться использовать git из командной строки.

\begin{enumerate}
    \item Ознакомиться с инструкцией по работе с git и GitHub из командной строки.
    \item Клонировать репозиторий командой \texttt{git clone <ссылка>}.
    \item Создать новый файл \texttt{info.md} в локальной копии репозитория.
    \item Записать в файл ФИО, название ОС, текущие дату и время~--- каждое с новой строки (3 строки).
    \item Выполнить \texttt{git add info.md} и \texttt{git commit}, ввести краткое описание коммита.
    \item Выполнить \texttt{git push}.
    \item Создать папку \texttt{data} в корне репозитория и файл \texttt{dataset.csv} в ней.
    \item Создать файл \texttt{.gitignore} со строкой \texttt{data/}.
    \item Закоммитить \texttt{.gitignore}, попробовать закоммитить \texttt{data/dataset.csv}, отправить коммиты через \texttt{git push}.
    \item Объяснить полученный результат в отчёте.
\end{enumerate}

\subsection*{Задание 3}
\textbf{Цель:} соблюсти формальности.

Подготовить отчёт по лабораторной работе (титульный лист, текст задания, результат выполнения). Сохранить отчёт в файле \texttt{report.pdf} в репозитории на GitHub.

\newpage
\section{Результаты выполнения}

\subsection{Задание 1. Файл \texttt{goals.md}}

\begin{code}
# Мой опыт
К сожалению, я не применял свои навыки программирования нигде, за пределами прохождения курсов ИТМО. Мой гитхаб: https://github.com/frogramer. Однако он заброшен, т.к. после прохождения дисциплины "Программирование" (Java) я ни разу так и не пользовался им
# Мои цели
1) Узнать, как обучают современные нейронные сети
2) Получить опыт в обучении нейронных сетей
3) Научиться работать с ИИ "по другую сторону" (со стороны разработчика)
## Задачи
- Научиться работать с GitHub
- Научиться обучать нейронные модели различными способами
- Научиться проверять обученные модели
![кот](resources/cat.jpg)
\end{code}

Изображение, добавленное в файл, находится в папке \texttt{resources} и вставлено в \texttt{goals.md} следующей строкой:

\begin{code}
![кот](resources/cat.jpg)
\end{code}

\subsection{Задание 2. Файл \texttt{info.md}}

Шаблон содержимого файла:

\begin{code}
<Фамилия Имя Отчество>
<Операционная система>
<ГГГГ-ММ-ДД ЧЧ:ММ>
\end{code}

Фактическое содержимое:

\begin{code}
Бузырев Иван Михайлович
Windows 11
2026-09-14 12:49
\end{code}

Коммит и отправка выполнены командами:

\begin{code}
git add info.md
git commit
git push
\end{code}

\subsection{Задание 2. Файл \texttt{.gitignore}}

Содержимое файла:

\begin{code}
data/
\end{code}

\subsection{Объяснение результата попытки закоммитить 

После создания папки \texttt{data} и файла \texttt{data/dataset.csv} был создан файл \texttt{.gitignore} со строкой \texttt{data/}, которая указывает Git игнорировать всё содержимое папки \texttt{data}. Файл \texttt{.gitignore} был успешно закоммичен командой:

\begin{code}
git add .gitignore
git commit
git push
\end{code}

Затем была предпринята попытка закоммитить файл \texttt{data/dataset.csv}:

\begin{code}
git add data/dataset.csv
git status
git commit
\end{code}

Git выдал предупреждение о том, что путь игнорируется одним из правил \texttt{.gitignore}:

\begin{code}
The following paths are ignored by one of your .gitignore files:
data
hint: Use -f if you really want to add them.
\end{code}

При выполнении \texttt{git commit} было получено сообщение \texttt{nothing to commit, working tree clean}~--- коммит не был создан, так как файл не попал в индекс. После \texttt{git push} папка \texttt{data/} и файл \texttt{dataset.csv} в удалённом репозитории на GitHub не появились.

\textbf{Вывод.} Файл \texttt{.gitignore} корректно исключает указанные пути из отслеживания Git. Это позволяет не загружать в репозиторий большие датасеты, временные файлы, кэш и другие данные, которые не должны версионироваться. При необходимости файл можно добавить принудительно командой \texttt{git add -f <путь>}, однако в рамках задания этого делать не требовалось. Следует также помнить, что \texttt{.gitignore} не влияет на уже отслеживаемые файлы: если файл был закоммичен до добавления правила, его необходимо сначала убрать из индекса командой \texttt{git rm {-}{-}cached <путь>}.

\section{Выводы}

В ходе лабораторной работы были освоены базовые приёмы работы с системой контроля версий Git и сервисом GitHub: создание файлов через веб-интерфейс, клонирование репозитория, работа с индексом, коммиты и отправка изменений через \texttt{git push}. Также было изучено назначение файла \texttt{.gitignore} и продемонстрировано, что пути, указанные в нём, не попадают в репозиторий.

\end{document}